\documentclass[a4paper,12pt]{article}
\usepackage[T2A]{fontenc}
\usepackage[utf8]{inputenc}
\usepackage[russian]{babel}
\usepackage{amsmath,amssymb,mathtools}
\usepackage{geometry}
\geometry{left=24mm,right=22mm,top=25mm,bottom=24mm}
\usepackage{PTSans}
\renewcommand{\familydefault}{\sfdefault}
\usepackage{microtype}
\usepackage{array,booktabs,tabularx}
\usepackage{enumitem}
\usepackage{tikz}
\usepackage{xcolor}
\usepackage{hyperref}
\usepackage{fancyhdr}
\usepackage{setspace}
\onehalfspacing
\definecolor{accent}{HTML}{087E82}
\definecolor{ink}{HTML}{29333C}
\definecolor{muted}{HTML}{647079}
\definecolor{err}{HTML}{AE3442}
\definecolor{idea}{HTML}{215B8F}
\hypersetup{colorlinks=true,linkcolor=accent,pdfauthor={Лёвин Артём Александрович},pdftitle={Ревью домашней работы — Ярослав Гаврилов — 02.11.2026}}
\pagestyle{fancy}
\fancyhf{}
\fancyhead[L]{\small\color{muted}Ревью домашней работы}
\fancyhead[R]{\small\color{muted}02.11.2026}
\fancyfoot[C]{\small\color{muted}\thepage}
\renewcommand{\headrulewidth}{0pt}
\setlength{\headheight}{15pt}
\setlength{\parindent}{0pt}
\setlength{\parskip}{5pt}
\setlist[itemize]{leftmargin=*,itemsep=2pt,topsep=2pt}
\newcommand{\sectiontitle}[1]{\vspace{6mm}{\large\bfseries\color{accent}#1}\par\vspace{2mm}}
\newcommand{\partlabel}[1]{\par\smallskip{\bfseries\color{ink}#1}\par\smallskip}
\begin{document}
\thispagestyle{empty}
\vspace*{38mm}
\begin{center}
{\small\bfseries\color{accent}МАТЕМАТИКА}\par
\vspace{19mm}
{\Huge\bfseries\color{ink}Ревью домашней работы}\par
\vspace{10mm}
{\Large Логарифмические уравнения и функции}\par
\vspace{27mm}
{\Large Ярослав Гаврилов}\par
\vspace{9mm}
{\large 2 ноября 2026 г.}\par
\end{center}
\vfill
\begin{center}
{\normalsize Лёвин Артём Александрович}\par
{\small\color{muted}эксклюзивно для Ярослава Гаврилова}\par
\vspace{12mm}
\begin{tikzpicture}
\draw[accent!65,line width=0.9pt] (-5,0) .. controls (-2.7,.65) and (-1.6,-.7) .. (0,0)
 .. controls (1.8,.65) and (3.8,-.6) .. (5,0);
\end{tikzpicture}
\end{center}
\newpage
\sectiontitle{Сводная таблица}
\small
\renewcommand{\arraystretch}{1.4}
\begin{tabularx}{\linewidth}{@{}p{12mm}p{29mm}X p{24mm}p{25mm}@{}}
\toprule
№ & Тема & Суть ошибки & Ваш ответ & Правильный ответ \\
\midrule
\hyperref[task:3]{3} & Логарифмы & Неверный знак после деления на отрицательное число; итог верен. & $x=0{,}2$ & $x=\dfrac15$ \\
\addlinespace[2mm]
\hyperref[task:9]{9} & Функция & Потерян минус перед свободным членом при вычислении $f(27)$. & $f(27)=1$ & $f(27)=1$ \\
\bottomrule
\end{tabularx}
\par\normalsize
\vspace{8mm}
Обратите внимание: оба итоговых ответа верны. В таблицу включены задания, где встречаются неверные промежуточные записи. Они могут привести к ошибке в похожей задаче.
\newpage
\thispagestyle{empty}
\vspace*{0.37\textheight}
\begin{center}
{\Huge\bfseries\color{accent}Разбор ошибок}
\end{center}
\newpage
\phantomsection\label{task:3}
\sectiontitle{Задание № 3. Сравнение логарифмов}
\partlabel{Текст задания}
Решите уравнение:
\[
\log_3(x+4)=\log_3(2x+1)+1.
\]
\partlabel{Ваше решение}
Вы верно представили единицу как $\log_3 3$, объединили логарифмы и получили
\[
\log_3(x+4)=\log_3(6x+3),\qquad
x+4=6x+3,\qquad -5x=-1.
\]
Затем в тетради записано $x=-\dfrac15$, а следующей строкой --- $x=+0{,}2$.
\partlabel{Ваш ответ}
$\boldsymbol{x=0{,}2}$ (верный).
\partlabel{Ошибка}
\textcolor{err}{\textbf{Ошибка:}} неверно определён знак частного при делении двух отрицательных чисел.
\partlabel{Где именно возникает ошибка}
Последний правильный шаг:
\[
-5x=-1.
\]
Первый неправильный шаг:
\[
\textcolor{err}{x=-\frac15}.
\]
\partlabel{Почему это неверно}
Чтобы найти $x$, обе части равенства $-5x=-1$ требуется разделить на $-5$:
\[
x=\frac{-1}{-5}.
\]
Частное двух отрицательных чисел положительно. Значит, $x=\dfrac15$. Отрицательный знак перед дробью в промежуточной строке противоречит правилу знаков. Десятичная запись $+0{,}2$ в следующей строке уже верна, но обе записи должны обозначать одно и то же число.
\partlabel{Ключевая идея}
\textcolor{idea}{\textbf{Ключевая идея:}} при делении обеих частей уравнения на отрицательное число знак результата определяется знаками числителя и знаменателя. Перед завершением решения полезно проверить равенство последовательно по всем строкам.
\partlabel{Исправление от последнего правильного шага}
\[
-5x=-1\quad\Longrightarrow\quad
x=\frac{-1}{-5}=\frac15=0{,}2.
\]
\partlabel{Полное правильное решение}
Сначала определим область допустимых значений. Аргументы обоих логарифмов должны быть положительными:
\[
x+4>0,\qquad 2x+1>0.
\]
Отсюда $x>-4$ и $x>-\dfrac12$. Совместное ограничение: $x>-\dfrac12$.

Применим равенства $1=\log_3 3$ и $\log_a u+\log_a v=\log_a(uv)$ при положительных $u$ и $v$:
\begin{align*}
\log_3(x+4)&=\log_3(2x+1)+\log_3 3,\\
\log_3(x+4)&=\log_3\bigl(3(2x+1)\bigr),\\
\log_3(x+4)&=\log_3(6x+3).
\end{align*}
Поскольку основания логарифмов совпадают и равны $3>0$, $3\ne1$, их аргументы равны:
\begin{align*}
x+4&=6x+3,\\
4-3&=6x-x,\\
1&=5x,\\
x&=\frac15.
\end{align*}
\partlabel{Проверка}
Число $\dfrac15$ удовлетворяет условию $x>-\dfrac12$. Для него
\[
x+4=\frac{21}{5},\qquad 2x+1=\frac75,\qquad
\frac{21/5}{7/5}=3.
\]
Поэтому $\log_3\dfrac{21}{5}=\log_3\dfrac75+\log_3 3=\log_3\dfrac75+1$.
\partlabel{Итог}
\textcolor{accent}{\textbf{Итог:}} $\boxed{x=\dfrac15}$. Знак в промежуточном шаге необходимо исправить.
\partlabel{Задача на закрепление}
Решите уравнение $\log_4(x+7)=\log_4(3x+1)+1$.
\newpage
\phantomsection\label{task:9}
\sectiontitle{Задание № 9. Логарифмическая функция}
\partlabel{Текст задания}
График функции $f(x)=\log_a x+b$ проходит через точки $(1;-2)$ и $(9;0)$. Найдите $f(27)$.
\partlabel{Ваше решение}
Вы верно использовали первую точку для нахождения $b=-2$, затем вторую точку для нахождения основания $a=3$. При вычислении значения функции записано
\[
\textcolor{err}{f(27)=\log_3 27+2},
\]
а ниже указан верный ответ $f(27)=1$.
\partlabel{Ваш ответ}
$\boldsymbol{f(27)=1}$ (верный).
\partlabel{Ошибка}
\textcolor{err}{\textbf{Ошибка:}} при подстановке аргумента потерян минус у свободного члена $b=-2$.
\partlabel{Где именно возникает ошибка}
Последний правильный результат --- $a=3$, $b=-2$. Поэтому функция имеет вид $f(x)=\log_3 x-2$.

Первый неправильный шаг --- переход к $f(27)=\log_3 27+2$ вместо $\log_3 27-2$.
\partlabel{Почему это неверно}
Число $b$ --- постоянное слагаемое функции. После того как получено $b=-2$, этот знак должен сохраняться при подстановке любого положительного $x$.

Если бы функция имела вид $\log_3 x+2$, то в точке $x=1$ её значение было бы $\log_3 1+2=2$, что противоречит условию $f(1)=-2$. Кроме того, записанное Вами промежуточное выражение $\log_3 27+2$ равно $3+2=5$ и не даёт итогового ответа $1$.
\partlabel{Ключевая идея}
\textcolor{idea}{\textbf{Ключевая идея:}} после нахождения параметров подставьте их в общую формулу и один раз запишите окончательный вид функции. Все дальнейшие значения вычисляйте уже по этой записи.
\partlabel{Исправление от последнего правильного шага}
\[
a=3,\quad b=-2\quad\Longrightarrow\quad
f(27)=\log_3 27-2=3-2=1.
\]
\partlabel{Полное правильное решение}
Для логарифмической функции необходимо $a>0$, $a\ne1$, $x>0$.

Точка $(1;-2)$ принадлежит графику. Поэтому
\[
-2=\log_a 1+b=0+b,\qquad b=-2.
\]
Точка $(9;0)$ тоже принадлежит графику:
\begin{align*}
0&=\log_a 9+b,\\
0&=\log_a 9-2,\\
\log_a 9&=2,\\
a^2&=9.
\end{align*}
Из $a^2=9$ получаем $a=3$ или $a=-3$. Основание логарифма обязано быть положительным, поэтому $a=3$.

Следовательно,
\[
f(x)=\log_3 x-2.
\]
Так как $27=3^3$,
\[
f(27)=\log_3 3^3-2=3-2=1.
\]
\partlabel{Проверка}
Для найденной функции
\[
f(1)=\log_3 1-2=-2,\qquad
f(9)=\log_3 9-2=2-2=0.
\]
Обе заданные точки действительно принадлежат графику; аргумент $27>0$ допустим.
\partlabel{Итог}
\textcolor{accent}{\textbf{Итог:}} $\boxed{f(27)=1}$. В промежуточной строке требуется сохранить $b=-2$.
\partlabel{Задача на закрепление}
График функции $g(x)=\log_a x+b$ проходит через точки $(1;-1)$ и $(16;1)$. Найдите $g(64)$.
\vfill
\begin{center}
{\footnotesize\color{muted}Лёвин Артём Александрович эксклюзивно для Ярослава Гаврилова}
\end{center}
\end{document}
